\subsection{Our Contributions}\label{subsec:contributions}

This paper provides a focused empirical and implementation oriented study of adaptive thresholding techniques applied to noisy PPM and JPEG color images.
Rather than proposing new denoising algorithms, we emphasize on the efficient realization and practical evaluation of established methods.

Our contributions lie in the careful adaption of color image data, the use of integral image representations to eliminate redundant computations, and a systematic performance analysis guided by runtime profiling.
Through comparative evaluation, we demonstrate how algorithmic choices affect denoising quality and how the implementation decisions impact the computational efficiency. 


% - implemented different approaches
% - compared different approaches
% - integral
% - binarization
%   - 
% - show different approaches next to each other

% % Optimizations

% - running our applications on Mac
% - using the "Instruments" Profiler inherent on mac
% - idea: only optimize what is necessary / causing a great bottleneck
% - baseline image (binarization)
%   - shows memory access of the two vectors for $windows_{mean}$ and $windows_{stddev}$
% - rewriting the loops (not calculate everything twice) (before $4^n$ not $2^n$) - using integral images
